\documentclass[uplatex,dvipdfmx,a4paper,11pt]{jsarticle}

\usepackage[top=25mm,bottom=25mm,left=25mm,right=25mm]{geometry}

\usepackage{enumitem}

\setlength{\parindent}{0pt}
\setlength{\parskip}{0.6em}

\usepackage{titlesec}
\titleformat{\section}
  {\normalfont\Large\bfseries}{\thesection.}{0.6em}{}
\titleformat{\subsection}
  {\normalfont\large\bfseries}{\thesubsection}{0.6em}{}
\titlespacing*{\section}{0pt}{1.2em}{0.6em}
\titlespacing*{\subsection}{0pt}{1.0em}{0.4em}

\renewcommand{\labelitemi}{\textbullet}
\renewcommand{\labelitemii}{\textopenbullet}
\setlist[itemize]{leftmargin=1.6em,itemsep=0.1em,topsep=0.2em}

\providecommand{\textopenbullet}{\ensuremath{\circ}}


\begin{document}

\section{No.2231 LEDマトリクス制御：計画書との変更点}

\subsection{描画APIの変更}

計画書（3.5.8.3）では，フレームバッファの描画に
\texttt{matrix.renderBitmap()}を使用すると記述されている．
\texttt{renderBitmap()}は\texttt{uint8\_t[8][12]}形式の配列を引数に取る関数である．
実装では\texttt{uint32\_t[3]}形式にビットパックし，
\texttt{matrix.loadFrame()}で描画する方式に変更した．
これは，\texttt{loadFrame()}の方がフレームバッファのメモリ効率が高く，
96ビット（8行$\times$12列）を3つの32ビット整数で表現できるためである．

\subsection{関数の追加}

計画書の表3.33には\texttt{checkUDP}と\texttt{displayBPM}の2関数のみ記載されている．
実装では以下の3関数を追加した．

\begin{itemize}
  \item \texttt{clearDisplay()}: 演奏終了ヘッダ（'E'）受信時にLEDマトリクスを全消灯する関数．
  計画書の通信プロトコル（表3.29, 3.30）にはヘッダ'E'が定義されているが，
  Display側の関数一覧にはこの処理への言及がなかったため追加した．

  \item \texttt{drawDigit()}: 1桁の数字を独自フォント（\texttt{font3x5[]}）を参照して
  フレームバッファに描画する内部関数．
  \texttt{displayBPM()}から桁ごとに呼び出される．

  \item \texttt{displaySetup()}: \texttt{matrix.begin()}をDisplay.cpp側で呼ぶための
  初期化関数．\texttt{Arduino\_LED\_Matrix}の表示バッファはヘッダ内でstatic定義されており，
  includeする翻訳単位ごとに別実体となる．そのため，描画を行う\texttt{loadFrame()}と
  同一ファイルから\texttt{begin()}を呼び出す必要があり，本関数を追加した．
\end{itemize}

\subsection{モジュール構成の変更}

計画書の図3.16では\texttt{checkUDP()}はDisplay.ino側に配置されている．
実装ではDisplay.cppに移動した．
\texttt{checkUDP()}内で\texttt{currentBPM}（Display.cppの静的変数）を参照する必要があるためである．
これにより，Display.inoには\texttt{setup()}と\texttt{loop()}のみが残る構成となった．

\subsection{受信データの読み込み形式}

計画書（3.5.8.3）では「ペイロードを文字列として読み込む」と記述されている．
実装では\texttt{uint8\_t buffer[2]}としてバイナリ形式で読み取り，
\texttt{buffer[0]}をヘッダ文字（char），\texttt{buffer[1]}を段階番号（uint8\_t）として処理している．
通信プロトコル仕様（表3.29, 3.30）自体は固定長2バイトのバイナリ形式で定義されており，
実装側がプロトコル仕様に忠実な形式を採用した．

\subsection{静的IP設定の追加}

計画書の表3.24にはDisplay側の定数として\texttt{localIP}（観覧車デバイス自身の静的IPアドレス）が
定義されている．実装では\texttt{WiFi.config(localIP, gateway, subnet)}を
\texttt{WiFi.begin()}の前に呼び出すことで，
指揮デバイスの送信先IPアドレス（192.168.11.10）と一致する固定IPを設定した．

\end{document}
